\section{Participation, Surrender, and Permission Frontiers}
\label{sec:s13}
This section proves the additional contracting results and specifies the numerical target of their implementation. The finite state space, bounded rewards, eight-interval calendar, first-action classification, and frozen later menus are those of S.12. A surrender action is available only at eligible interior nodes. Its reward is $G-F$ and its live transition row is zero. Boundary discharge continues to pay $G$. The first interval remains compulsory in every reported position-class experiment.

\subsection{Affine Payoffs, Participation, and Procurement}
For a deterministic Markov policy, augment the original action labels by a surrender label. Add an absorbing cemetery state if needed to represent killed probability mass. On reaching the surrender label, pay $G-F$ once and receive zero thereafter. The unaltered terminal payoff is $G$. Backward substitution in this policy recursion, or direct summation over the finite path space, gives
\[
 J_p=\mathcal B_p+d\mathcal A_p-F\mathcal S_p.
\]
The operating duration is nonnegative and bounded by $(1-e^{-\rho})/\rho$. An elective surrender occurs at most once, so $0\leq\mathcal S_p\leq1$. In particular, the fee exposure cannot be replaced by the expected number of visits to surrender-eligible states.

There are finitely many deterministic Markov policies in the original action target. They attain the unrestricted and sign-restricted optima; randomization does not improve a finite-horizon expected payoff. Each policy payoff is affine in $(d,F)$, so their maximum is convex, nondecreasing in $d$, and nonincreasing in $F$. If $F'\geq F$, every fixed policy satisfies $-(F'-F)\leq J_p(F')-J_p(F)\leq0$. Taking policy maxima on the appropriate side gives the $1$-Lipschitz property. Increasing the minimum term removes surrender policies, so it weakly lowers each class value. None of these orders implies an order for a difference of class values.

For completeness, consider a direction $v=(v_d,v_F)$ at a contract $z=(d,F)$. A finite set of nonoptimizing policies has a strictly positive value gap from the maximum, unless it is empty. For sufficiently small positive steps, only policies optimizing at $z$ can determine the directional derivative. The derivative is therefore
\[
 D_vW(z)=\max_{p\in\argmax J(z)}\{v_d\mathcal A_p-v_F\mathcal S_p\}.
\]
This proves the exposed-feature statement, including the case of several policies with the same feature vector. It is the standard finite envelope logic used in the main paper, not a new general envelope theorem.

The external numeraire makes total agent utility $q-C(E)+\max_\sigma W^r_\sigma$. Comparing this quantity with $G(x_0)$ and imposing $q\geq0$ gives the stated participation grant. The managed initial state is unchanged, so the grant does not enter the Bellman recursion. If a principal receives flow $b$ while the mandate operates, its expected payoff at the minimum grant is $b\mathcal A_{p_r^*}-q_r^*(E)$. Termination charges go to the separate enforcement institution. Dividing by strictly positive duration gives the procurement condition. If an agent optimum is not unique, the contract must specify an admissible selection or use the relevant worst-case duration; the proposition uses a specified selection. It does not assume that all payoff-equivalent policies deliver identical service.

Both regime policy sets use the same financial opportunities and outside payoff. The no-adjustment policies are feasible in the adjustment regime. Therefore $W^{\mathrm{adj}}_\sigma\geq W^0_\sigma$ and $q^{\mathrm{adj}}_\sigma\leq q^0_\sigma$ at a common capacity cost. Where all class-specific participation gaps are positive, subtracting the two participation equations within each class gives $q^0_\sigma-q^{\mathrm{adj}}_\sigma=W^{\mathrm{adj}}_\sigma-W^0_\sigma$. The four-class contrast yields the price interpretation of the relative option. The sign of that contrast is not implied by its two component inequalities.

The separate-account assumption is economically material. A model in which all payments enter managed CRRA wealth would require a new state transition, consumption response, and participation calculation. Similarly, collateral opportunity costs depending on realized duration would be additional running or terminal rewards, rather than the ex ante $C(E)$ used here. These are alternative economic targets, not accounting details that can be inserted without reoptimization.

\subsection{The Statewise Enforcement Frontier}
Let $R_n$ be the live support reachable using the union of endpoint kernels and all original feasible actions. Nonnegative kernel entries imply support closure: if $x\in R_n$ and a live row from $x$ gives positive mass to $y$, then $y\in R_{n+1}$. Elective surrender can remove paths, but cannot introduce a new live state. The union may be larger than the support of any particular optimized policy; this is intentional.

Suppose $F\geq F_r^*(m;\lambda,d)$. At maturity the optional and compulsory values both equal $G$. Suppose they coincide at all reachable nodes at date $n+1$. By support closure their operating backups coincide at any reachable node at date $n$. At an eligible interior node, $G-F\leq V^r_n$ by the definition of $F_r^*$; taking the maximum with surrender cannot raise the value. At an ineligible node the operating backup is the only option, and at a boundary the payoffs agree by construction. Backward induction proves statewise equality, and then equality of both sign-restricted initial values. If instead $F<F_r^*$, some eligible reachable interior node has $G-F>V^r_n$, at which its optional value is strictly higher. This proves the necessary direction for equality at every reachable node. Such a node might not affect either initial optimum, which is why the threshold is not described as the smallest initial-sign-preserving fee.

Increasing $m$ removes terms from the maximum. Increasing $d$ weakly raises $V^r_n$ at every state by nonnegative policy duration, so it weakly reduces the maximum. Policy inclusion gives $V^{\mathrm{adj}}_n\geq V^0_n$, and hence $F^*_{\mathrm{adj}}\leq F^*_0$. Thus adjustment weakly reduces the capacity needed to implement the same mandatory operating values. When $C$ is nondecreasing, an available smaller implementing capacity can also weakly reduce its participation cost. This claim concerns implementation capacity, not monotonicity of the initial risk preference.

The regional calculation uses the two no-adjustment noncancellable policies optimized at $(\lambda,d)=(0,0.4)$ and $(0.25,0.4)$. Their later-date policies are feasible regardless of the initial class used to store the policy. For each policy and each date $n\geq1$, compute its Bernstein coefficients in the \emph{physical} mixture probability $\lambda$. Restrict those coefficients to each of sixteen subintervals of $[0,0.25]$. Their coefficient minimum is a uniform policy lower value at every state on that subinterval. Take the larger of the two policy lower values, and subtract $10^{-7}$ before comparing with $G$. Since duration is nonnegative, this lower bound also applies for all $d\in[0.4,0.45]$. It is below $V^0$ and therefore below both regime values. Maximizing the resulting liquidation deficit over eligible reachable nodes gives the sufficient regional frontier.

The resulting bound $0.8348592805606$ for $m=1$ includes the stated arithmetic allowance. The upward-rounded value $0.834860$ is sufficient relative to the stored arrays. The optimizing deficit in this bound occurs at date one, node $(1.95,1.34375)$, on the first law subinterval. Restricting to the actual initial state alone would not justify ignoring that continuation constraint. Conversely, including every unrelated grid node raises the sufficient bound to $2.5309784144$. Neither bound is a claim about a continuum of initially admissible managed states.

\begin{table}[tbp]
\centering\small
\caption{Statewise No-Surrender Thresholds at the Original Center}
\label{tab:r9_center_capacity}
\begin{tabular}{rrr}
\toprule
$m$ & Adjustment & No adjustment \\
\midrule
1 & 0.7767542 & 0.8135332 \\
2 & 0.7161762 & 0.7458682 \\
3 & 0.6442995 & 0.6661886 \\
4 & 0.5592377 & 0.5734029 \\
5 & 0.4575809 & 0.4648691 \\
6 & 0.3346276 & 0.3367796 \\
7 & 0.1842712 & 0.1842768 \\
8 & 0.0000000 & 0.0000000 \\
\bottomrule
\end{tabular}
\begin{minipage}{0.98\textwidth}\footnotesize\medskip
Direct evaluations of $\max_{n\geq m,\,x\in R_n}[G(x)-V_n^r(x)]_+$ at $(\lambda,d)=(0.125,0.425)$. Both use the same all-action reachable support. The initial point alone can require a smaller charge than statewise equality. These point calculations are distinct from the upward-rounded sufficient regional bounds in the main paper.
\end{minipage}
\end{table}

Table \ref{tab:r9_center_capacity} gives the direct-solver statewise thresholds at the original center. They use the same reachable support and the separately optimized mandatory continuations. In particular, the one-interval thresholds are approximately $0.7767542$ with adjustment and $0.8135332$ without adjustment. They illustrate the policy-inclusion order. The sufficient regional bounds in the main text additionally cover all contracts in the rectangle and use feasible-policy lower values; the two tables answer different questions.

If $G$ is an operating supersolution at every date, the zero-charge recursion has $W_n=\max(G,\mathcal T_nG)=G$ at the last eligible decision date. Repeating the argument proves equality at all eligible dates. Earlier mandatory dates retain their operating recursion. At date zero, admitting free surrender yields $G(x_0)$ as well. In the finite implementation, every operating action is checked at both law endpoints and the upper benefit endpoint. Since the one-step expression is affine in the law weight and its duration coefficient is nonnegative, these endpoint inequalities cover the entire original rectangle. The large observed slack is a supersolution check for the stored arrays, not an interval proof for an uncomputed transition constructor.

\subsection{Risky-Share Breakpoints and Permission Derivatives}
Fix an initial $(c,\theta)$ pair. For each law and sign realization the proposed preference coordinate is $2+h\theta+\sqrt h\sigma_u z_u$, independent of $\pi$. The wealth coordinate has the form
\begin{align*}
 y_X(\pi)&=a_c+b_{\varrho,z}\pi,\qquad a_c=1.25+h(r\,1.25-c),\\
 b_{\varrho,z}&=h(\mu-r)1.25+\sqrt h\sigma_X1.25
             (\varrho z_u+\sqrt{1-\varrho^2}z_X).
\end{align*}
For the declared constants, $|y_u-2|<0.043$ and $1.054<y_X<1.365$ throughout the wider risky-share domain, all common $(c,\theta)$ pairs, and both law endpoints. These conservative bounds lie strictly inside $[1.2,2.8]\times[0.5,2]$. The whole line segment therefore stays interior, the exit fraction is one, and the flow annuity is independent of $\pi$.

Conditional on a shock, the preference interpolation weights are constant as $\pi$ varies. Between any two wealth-node crossings, each wealth interpolation weight is affine in $y_X$, hence affine in $\pi$. The bilinear interpolant of any continuation vector is consequently affine in $\pi$ there. Averaging shocks and the law mixture preserves affinity on the common partition. The flow reward is independent of $\pi$, so the first-date action value is continuous and piecewise affine. A maximum over a compact interval is attained at an endpoint or a breakpoint. A positive open interval uses its closure only to compute the supremum; the implementation reports the chosen action and does not call a zero endpoint a positive maximizer. The original finite frozen proposals are compared separately when feasible.

For an interpolation cell, write $I_XV_1(y)$ for the slope of the interpolant in wealth. Differentiation gives
\[
 \partial_\pi Q(c,\theta,\pi)
 =\sum_{\varrho}\Pr(\varrho)\frac14\sum_z
    e^{-\rho h} I_XV_1(y_{\varrho,z}(\pi))\,b_{\varrho,z}.
\]
There is no within-period flow derivative in this expression. If the winning pair is locally fixed and the relevant optimum remains at $L$ or $-S$, the value of relaxing the permission is the corresponding endpoint derivative. At a breakpoint or an action tie, use the appropriate one-sided piecewise derivative or directional envelope. Because the first-date permissions do not change $V_1$, no later-menu derivative is being omitted.

Changing $d$, unlike changing only a first-date permission, changes the optimized continuation. The finite envelope theorem nevertheless gives the derivative of each class value as its optimizing discounted duration. Differentiating $\Delta^r(\lambda,d_r^*,F,L,S)=0$ at a regular crossing then gives the three frontier formulas in the main text. If $D_A^r=0$, if a feature vector is not exposed, or if a control switch destroys differentiability, these ratios are not asserted. The deposited benefit-root calculations reoptimize continuation at every trial benefit and retain local endpoint signs. Their floating-point root brackets are not directed-rounding interval enclosures and do not prove that no other crossing exists.

\subsection{The Joint Fee--Benefit--Law Certificate}
A fixed stopping policy evaluated under affine endpoint kernels has a polynomial value of degree at most the remaining horizon in $\lambda$. This follows directly from the recursion: each live step multiplies by one affine mixture, whereas a surrender row has zero live kernel. Padding with zero future rewards and a cemetery state permits degree elevation to the common horizon. Applying the same recursion to the base reward, duration reward, and unit surrender charge yields the three focal coefficient vectors $B_p,A_p,S_p$.

At each $(d,F)$ corner, the signed correction and the count recursion use the operating and surrender actions together. Their proofs in S.10 apply because the added action has a bounded affine reward and a nonnegative zero kernel. The last-decision endpoint continuations agree at the terminal vector, and the initial sign mask is applied only when forming the date-zero class upper values. The count information and the signed correction are not granted a different surrender date or a different charge from the feasible lower policies.

For a fixed law probability, the optimized value is a maximum of affine functions of $(d,F)$. Its value inside a rectangle is therefore bounded above by the bilinear convex combination of its values at the four corners. Replacing those corner values by valid law-polynomial uppers preserves the inequality. A fixed policy value is affine in $(d,F)$ and equals the corresponding corner interpolation. Restrict both the upper and policy polynomials to a law subinterval and express them in a common Bernstein degree. For a positive policy $p$ and nonpositive upper $U_-$, the minimum of all coefficients of $J_p-U_-$, including the reward-parameter corner indices, bounds that difference everywhere in the cell. Taking the largest such bound across positive policies remains valid. The upper difference is obtained symmetrically from $U_+-J_q$ and a minimum across nonpositive policies. Taking the worst cell gives the regional statement. No policy is selected using information unavailable in its actual evaluation.

The deposit constructs the full four-corner fee--benefit bank at the two law endpoints for both adjustment regimes and both initial classes. The $F\in[0.85,0.90]$ certificate uses four law subcells, while retaining the original full menu plus the eligible surrender action. The favorable signs are not inferred by extrapolating the pointwise no-surrender observation at the center. The reachable enforcement bound supplies an additional, separate sufficient explanation for their equality with the noncancellable signs. The $F=0.80$ certificate fails; direct reoptimization at $(\lambda,d)=(0,0.4)$ gives a positive no-adjustment difference $0.0013923796303$, so this failure is not merely an uninformative upper bound.

\subsection{Arithmetic, Target Transfer, and Executed Validation}
The arithmetic assumptions are those in S.12: stored reward arrays and nonnegative transition arrays are evaluated in round-to-nearest binary64 arithmetic. Rewards can have either sign; nonnegativity applies to transition weights and duration, not to utility. The surrender reward is included in the absolute reward bound and its row mass is zero. Let $\gamma_j=j2^{-53}/(1-j2^{-53})$. The R9 audit reuses the $B,E,M,D$ backward recursions in S.12 with row-mass bound $\beta=0.9950124791936824$, absolute reward bound $R=27.2130993863$, and law width $1/4$. The bound covers $|d|\leq1$ and $0\leq F\leq1$.

For the additional three-feature lower representation, the audit uses
\[
 L_{\rm err}=5E_0+\gamma_{4096(8+1)}(1+8B_0),
\]
and the same upper error as in S.12, enlarged finally by
$\gamma_{64}(1+4B_0+2M_0)$ for coefficient subtraction and extrema. The factor five charges the subtraction of the base value in constructing the duration and surrender features and their recombination, instead of treating those features as exact. The restriction budget dominates both de Casteljau passes. A sparse row has at most sixteen entries before aggregation; adding one surrender label does not increase a dot-product length. Maxima are nonexpansive and introduce no factor equal to the number of feasible actions. The resulting derived per-class bound is $1.4305784\times10^{-8}$, below the deliberately larger accepted allowance $10^{-7}$. Every displayed difference interval is enlarged by $2\times10^{-7}$.

The stored-array target is not silently identified with a real-arithmetic transition constructor. For two finite targets with a common action correspondence, subtract their Bellman backups, add and subtract one kernel applied to the other's continuation, and use
$|\max_a q_a-\max_a\widetilde q_a|\leq\max_a|q_a-\widetilde q_a|$.
An absolute reward discrepancy $\epsilon_r$, row-$\ell^1$ discrepancy $\epsilon_K$, continuation magnitude bound $B$, and row-mass bound $\beta$ give
$\epsilon_n\leq\epsilon_{r,n}+\epsilon_{K,n}B_{n+1}+\beta_n\epsilon_{n+1}$.
The terminal discrepancy initializes this recursion. Surrender, forced-discharge, and boundary payments must enter the reward or terminal discrepancy. Different action sets require a separate action-approximation bound or a correspondence with a priced improvement error. A diffusion additionally requires its own consistency and approximation argument. The arithmetic allowance already charged to the stored-array certificate is not counted a second time as a measured constructor error.

The new independent validation has two components. Twelve random two-state, three-interval models are checked by enumeration of 81 deterministic continuation policies per initial class and parameter point. There are 120 class--point comparisons. The largest discrepancy between enumeration and the production stopping recursion is $\RnineToyReplay$; count never exceeds the signed chord beyond the declared numerical tolerance. A separate full-economy check evaluates 80 class--parameter comparisons across four fees, both regimes, and five mixture probabilities. Selected-control reconstruction without the stored sparse rewards and kernels differs from the optimized value by at most $\RnineFullReplay$. Fixed-policy polynomial replay differs by at most $\RninePolicyReplay$. Additional checks cover reward-parameter convexity, action eligibility, and the distinct date-zero free-surrender contract. These are implementation checks, not substitutes for the preceding proofs.

The complete economic run retains 19 center surrender/term specifications, 20 further fee--contract rows, their four class policies and moments, sixteen all-action supersolution checks, the eight sufficient enforcement bounds, nine two-regime permission comparisons, and eighteen reoptimized benefit roots. Three regional fee specifications are checked with both streamed upper constructions. All 398 protected historical or numerical source entries checked by the execution manifest remain unchanged. The permission target is the continuous first-date risky-share envelope conditional on finite consumption--adjustment pairs; the fee certificate target is the original finite menu plus surrender. The deposit does not claim a joint theorem simultaneously enlarging those two different action targets.

\subsection{Reproduction}
The complete manuscript entry points are \texttt{ECTA\_R9.tex} and \texttt{SUPP\_R9.tex}. From the repository root, run the following commands in order:
\begin{quote}\small\ttfamily
python replication/r9/run.py --mode all\\
python replication/r9/validate.py\\
python replication/r9/tables.py
\end{quote} The source inventory, historical hash checks, generated JSON, policy and coefficient arrays, execution transcript, manuscript build log, and point-by-point response are deposited together. R9 adds new readable source files and new evidence directories. The original R8 manuscript, supplement, all earlier economic applications and adverse comparisons, and the latest advisory report remain available at their original paths.
